\documentclass[letterpaper]{article}

\usepackage[header]{aimatters}


% Complete the title and author information below
\title{Formatting Instructions for Authors Submitting to AI Matters}
\author{
\AIMauthor{AI Matters Editors}{ACM SIGAI}{aimatters@sigai.acm.org}  % Fill in author information, one per line
\AIMauthor{Shakey the Robot}{Stanford Research Institute}{shakey@sri.com}
}


% Add a bio for each author.  (These will appear at the end of the document.)
% Note:  Due to quirks in the way LaTeX wraps text around figures, only use \textbf or \textem for emphasis in the bios; no other formatting commands!  (Specifically, don't use \bf, \em, etc.)
\AIMbio{sigai.jpg}{\textbf{The AI Matters Editors} may be contacted at \href{mailto:aimatters@sigai.acm.org}{aimatters@sigai.acm.org} with questions about the submission process for the AI Matters newsletter.  Complete submission instructions may be found on the AI Matters website.  If you note any bugs in the AI Matters template, please let the editors know.}

\AIMbio{shakey.jpg}{\textbf{Shakey the robot} was the first general-purpose mobile robot to be able to reason about its own actions. The project combined research in robotics, computer vision, and natural language processing. Because of this, it was the first project that melded logical reasoning and physical action. Shakey was developed at the Artificial Intelligence Center of Stanford Research Institute (now called SRI International).}


\begin{document}
\maketitle
 

\section{Introduction}

This document details the requirements to prepare your article for publication in AI Matters.  All submissions to AI Matters must be formatted in \LaTeX{} using this template.  These instructions were adapted from the AAAI formatting instructions.

These instructions are provided as a general guide for experienced \LaTeX{} users.  AI Matters staff cannot provide you with support in formatting your submission, and provide no warranties or guarantees that the accompanying {\em aimatters.sty} style files will work with all systems.  However, if you notice any errors in the provided style file, please email the current AI Matters editors at \href{mailto:aimatters@sigai.acm.org}{aimatters@sigai.acm.org} to let us know.  If the results you obtain are not in accordance with the specifications you received, you must correct your source file to achieve the correct result; please do not edit the provided {\em aimatters.sty} file, since we will be using a fresh copy of the style file to compile your final submission. 

All authors must comply with the following:

\begin{itemize}
\item You must use the latest AI Matters \LaTeX{} style file {\em aimatters.sty}, which can be downloaded from the AI Matters website.
\item Read and format your paper source and PDF according to the formatting instructions for authors.
\item Submit your \LaTeX{} source files, compiled PDF, and all necessary files (images, .bib files, etc.) to re-build your submission.
\item Follow the submission instructions posted on the AI Matters website.
\item Check every page of your paper before submitting it.
\end{itemize}

When you build your PDF using the AI Matters template, you will notice a large black box next to the title.  Do not worry about this; the editors will replace it with the article category when compiling the final version of the newsletter.  Similarly, the DOI for your article and newsletter volume/number/date will be all X's as placeholders.  Again, the editors will take care of filling in the correct DOI for your article, and the newsletter volume/number/date -- you do not need to do anything with these.


\section{Copyright \& Notice to Contributing Authors to SIG Newsletters}
By submitting your article for distribution in this Special Interest Group publication, you hereby grant to ACM the following non-exclusive, perpetual, worldwide rights:
\begin{itemize}
\item to publish in print on condition of acceptance by the editor
\item to digitize and post your article in the electronic version of this publication
\item to include the article in the ACM Digital Library and in any Digital Library related services 
\item to allow users to make a personal copy of the article for noncommercial, educational or research purposes.
\end{itemize}
However, as a contributing author, you retain copyright to your article and ACM will refer requests for republication directly to you.


\section{Formatting Requirements in Brief}
We need source and PDF files that can be used in a variety of ways and can be output on a variety of devices. AI Matters imposes some requirements on your source and PDF files that must be followed by all submissions:
\begin{itemize}
\item Your .tex file must compile in PDF\LaTeX{} --- \textbf{ no .ps or .eps figure files.}
\item All fonts must be embedded in the PDF file --- \textbf{ this includes your figures.}
\item Modifications to the style sheet are NOT allowed.
\item No type 3 fonts may be used (even in illustrations).
\item Your title must follow US capitalization rules.
\item No \LaTeX{} 209 documents may be used or submitted.
\item Fonts that require non-English language support (CID and Identity-H) must be converted to outlines or removed from the document (even if they are in a graphics file embedded in the document). 
\item Two-column format as specified in the style file is required for all papers.
\item The paper size for final submission must be US letter. No exceptions.
\item The source file must exactly match the PDF.
\item The document margins must be as specified in the formatting instructions.
\item No document may be password protected.
\item Your source and PDF must not have any page numbers, footers, or headers (besides those created by the style file).
\item Your PDF must be compatible with Acrobat 5 or higher.
\end{itemize}


\section{What Files to Submit}
You must submit the following items to ensure that your paper is published:
\begin{itemize}
\item A fully-compliant PDF file.
\item Your  \LaTeX{}  source file submitted as a \textbf{single} .tex file (do not use the ``input" command to include sections of your paper --- every section must be in the single source file). The only exception is the bibliography, which you may include separately. Your source must compile on our system using PDF\LaTeX{}, which includes the standard \LaTeX{} support files.
\item All your graphics files.
\item The \LaTeX{}-generated files (e.g. .aux and .bib file, etc.) for your compiled source.
\item All the nonstandard style files (ones not commonly found in standard \LaTeX{} installations) used in your document (including, for example, old algorithm style files). If in doubt, please include it.
\end{itemize}

Your \LaTeX{} source will be reviewed and recompiled on our system.   \textbf{Please do not submit your source in multiple text files.} Your single \LaTeX{} source file must include all your text, your bibliography (formatted using apacite.bst), and any custom macros. Accompanying this source file, you must also supply any nonstandard (or older) referenced style files and all your referenced graphics files. 

Your files should work without any supporting files (other than the program itself) on any computer with a standard \LaTeX{} distribution. 

\subsection{Naming your Files}

We recommend that you place your PDF and source files in a single tar, zipped, gzipped, stuffed, or compressed archive. Name your source file with the first author's last (family) name, so that it can easily be differentiated from other submissions.  Name the compressed archive of your submission similarly.


\section{Using \LaTeX{} to Format Your Paper}

The latest version of the {\em aimatters.sty} style file is available on the AI Matters website. Download this file and place it in a file named ``aimatters.sty" in the \TeX\ search path. Placing it in the same directory as the paper should also work. You must download the latest version of the style file and formatting instructions.

\subsection{Document Structure}

In the \LaTeX{} source for your paper, you \textbf{must} place the following lines as shown in the example in this subsection. This example is set up for three authors. Add or subtract author and address lines as necessary. Please do not change the default fonts; in order to ensure this, be sure to load any other packages before including the AI Matters style via the ``\textbackslash usepackage\{aimatters\}'' line.

\begin{quote}
\begin{tiny}
\textbackslash documentclass[letterpaper]\{article\}\\
\textbackslash usepackage[header]\{aimatters\}\\
\\
\% Complete the title and author information below\\
\textbackslash title\{Insert Your Title Here\}\\
\textbackslash author\{\\
\textbackslash AIMauthor\{Name1\}\{InstitutionA\}\{name1@instA.edu\}\\
\textbackslash AIMauthor\{Name2\}\{InstitutionA\}\{name2@instA.edu\}\\
\textbackslash AIMauthor\{Name3\}\{InstitutionB\}\{name3@instB.com\}\\
\}\\
\% Complete bios for each author.\\
\% (These will appear at the end of the document.)\\
\% Note:  Use only \textbackslash textbf or \textbackslash textem for emphasis in the bios;\\
\% no other formatting commands!  (E.g., don't use \textbackslash bf, \textbackslash em, etc.)\\
\textbackslash AIMbio\{author1.jpg\}\{\textbackslash textbf\{Author1\} is a PhD student at....\}\\
\textbackslash AIMbio\{author2.jpg\}\{\textbackslash textbf\{Author2\} is a professor....\}\\
\textbackslash AIMbio\{author3.jpg\}\{\textbackslash textbf\{Author3\} works at....\}\\
\\
\textbackslash begin\{document\}\\
\textbackslash maketitle\\
\\
... \% \textit{Put body of paper here}\\
\\
\% \textit{References and End of Paper}\\
\textbackslash bibliography\{Bibliography-File\}\\
\textbackslash bibliographystyle\{apacite\}\\
\textbackslash end\{document\}
\end{tiny}
\end{quote}



\subsection{Layout Requirements}
Papers must be formatted to print in two-column format on 8.5 x 11 inch US letter-sized paper, as specified in the AI Matters style file.  Do not use the Geometry package to alter the page size.  Text should not extrude into the column separation or the margins.

Commands that alter page layout are forbidden. These include \textbackslash columnsep, \textbackslash topmargin, \textbackslash topskip, \textbackslash textheight, \textbackslash textwidth, \textbackslash oddsidemargin, and \textbackslash evensizemargin (this list is not exhaustive).


\subsection{Figures}
Your paper must compile in PDF\LaTeX{}. Consequently, all your figures must be .jpg, .png, or .pdf. You may not use the .gif (the resolution is too low), .ps, or .eps file format for your figures.

When you include your figures, you must crop them \textbf{outside} of \LaTeX{}. The command \textbackslash includegraphics*[clip=true, viewport 0 0 10 10]{...} might result in a PDF that looks great, but the image is \textbf{not really cropped.} The full image can reappear when page numbers are applied or color space is standardized. 

\subsection{Type Font and Size}
Your paper must be formatted in Helvetica.  Please do not change the default font or default font sizes.  Use Symbol or Lucida or Computer Modern for \textit{mathematics only. } 

Do not use type 3 fonts for any portion of your paper, including graphics. Type 3 bitmapped fonts are designed for fixed resolution printers. Most print at 300 dpi even if the printer resolution is 1200 dpi or higher. They also often cause high resolution imagesetter devices and our PDF indexing software to crash. Consequently, AI Matters will not accept electronic files containing obsolete type 3 fonts. Files containing those fonts (even in graphics) will be rejected. 

If you are unsure if your paper contains type 3 fonts, view the PDF in Acrobat Reader. The Properties/Fonts window will display the font name, font type, and encoding properties of all the fonts in the document. If you are unsure if your graphics contain type 3 fonts (and they are PostScript or encapsulated PostScript documents), create PDF versions of them, and consult the properties window in Acrobat Reader. 

\subsubsection{Obtaining Type 1 Computer Modern for \LaTeX{}.}

If you use Computer Modern for the mathematics in your paper (you cannot use it for the text) you may need to download type 1 Computer fonts. They are available without charge from the American Mathematical Society:
http://www.ams.org/tex/type1-fonts.html. 

\subsection{Title and Authors}
Your title must appear in mixed case (nouns, pronouns, and verbs are capitalized) at the top of the first page.  Authors' names should appear below the title of the paper, along with affiliation(s) and email address(es).  We have included the \textbackslash AIMauthor macro to help format the authors lines.  It accepts three arguments: the name, the affiliation, and the email address, and will produce the correct formatting automatically.  For example, with three authors:

\begin{quote}
\begin{tiny}
\% Complete the title and author information below\\
\textbackslash title\{Insert Your Title Here\}\\
\textbackslash author\{\\
\textbackslash AIMauthor\{Sally Smith\}\{University of XYZ\}\{smith@xyz.edu\}\\
\textbackslash AIMauthor\{John Doe\}\{ABC Corporation\}\{doe@abc.com\}\\
\textbackslash AIMauthor\{Zoey Johns\}\{IJK College\}\{zjb7@ijk.edu\}\\
\}\\
\end{tiny}
\end{quote}



\subsection{Credits}
Any credits to a sponsoring agency should appear in the acknowledgments section, unless the agency requires different placement. If it is necessary to include this information on the front page, use
\textbackslash thanks in either the \textbackslash author or \textbackslash title commands.
For example:
\begin{quote}
\begin{small}
\textbackslash title\{Very Important Results in AI\textbackslash thanks\{This work is
 supported by everybody.\}\}
\end{small}
\end{quote}
Multiple \textbackslash thanks commands can be given. Each will result in a separate footnote indication in the author or title with the corresponding text at the botton of the first column of the document. Note that the \textbackslash thanks command is fragile. You will need to use \textbackslash protect.

Please do not include \textbackslash pubnote commands in your document.

\subsection{Abstract}
If included, the abstract must be placed at the beginning of the first column using the \LaTeX{} abstract environment.  This concise, one-paragraph summary should describe the general thesis and conclusion of your paper. A reader should be able to learn the purpose of the paper and the reason for its importance from the abstract. The abstract should be no more than two hundred words in length. (Authors who are submitting short one- or two-page extended extracts should provide a short abstract of only a sentence or so.) Do not include references in your abstract.

\subsection{Page Numbers}

Do not print any page numbers manually on your paper.  The AI Matters style file will add these automatically.

\subsection{Text }
The main body of the paper must be formatted in eleven-point with twelve-point leading (line spacing).

\subsection{Citations}
Citations within the text should use APA format.

\subsection{Footnotes}
Avoid footnotes as much as possible; they interrupt the reading of the text. When essential, they should be consecutively numbered throughout with superscript Arabic numbers. Footnotes should appear at the bottom of the page, separated from the text by a blank line space and a thin, half-point rule. 

\subsection{Headings and Sections}
When necessary, headings should be used to separate major sections of your paper.   AI Matters does not use section numbers.

\subsubsection{Acknowledgments.}
The acknowledgments section, if included, appears after the main body of text and is headed ``Acknowledgments." This section includes acknowledgments of help from associates and colleagues, credits to sponsoring agencies, financial support, and permission to publish. Please acknowledge other contributors, grant support, and so forth, in this section. Do not put acknowledgments in a footnote on the first page. If your grant agency requires acknowledgment of the grant on page 1, limit the footnote to the required statement, and put the remaining acknowledgments at the back. Please try to limit acknowledgments to no more than three sentences. 

\subsubsection{Appendices.}
Any appendices follow the acknowledgments, if included, or after the main body of text if no acknowledgments appear. 

\subsubsection{References}
The references section should be labeled ``References" and should appear as the second-to-last section of the paper (don't end the paper with references, and then put a figure by itself on the last page). A sample list of references is given later on in these instructions. Please use APA format for references. Please prepare complete and accurate citations.

\LaTeX{} and BibTeX support APA-style references automatically through the {\em apacite} bibliography style.

\subsubsection{Author Biographies}

The last section of your paper should be the author biographies.  A biography and photo is required of every author.  To insert each author's photo and bio, use the provided \textbackslash AIMbio command, which takes in two arguments:  an image and the biography text.

For example, 
\begin{quote}
\begin{tiny}
\textbackslash AIMbio\{shakey.jpg\}\{\textbackslash textbf\{Shakey the robot\} was the first general-purpose mobile robot....\}\\[3pt]
\textbackslash AIMbio\{data.jpg\}\{Lieutenant Commander \textbackslash textbf\{Data\} is a character in the fictional Star Trek universe....\}
\end{tiny}
\end{quote}

Due to quirks in the way LaTeX wraps text around figures, only use \textbackslash textbf or \textbackslash textem for emphasis in the bios; no other formatting commands!  (Specifically, don't use \textbackslash bf, \textbackslash em, etc.)

You may place these \textbackslash AIMbio commands anywhere in the document, and they will be automatically added to the end of the article.


\subsection{Illustrations and Figures}
Figures, drawings, tables, and photographs should be placed throughout the paper near the place where they are first discussed. Do not group them together at the end of the paper. If placed at the top or bottom of the paper, illustrations may run across both columns. Figures must not invade the top, bottom, or side margin areas. Figures must be inserted using the \textbackslash usepackage\{graphicx\}. Number figures sequentially, for example, figure 1, and so on. Figures should have appropriate captions.

\subsubsection{Low-Resolution Bitmaps.}
You may not use low-resolution (such as 72 dpi) screen-dumps and GIF files---these files contain so few pixels that they are always blurry, and illegible when printed. If they are color, they will become an indecipherable mess when converted to black and white. This is always the case with gif files, which should never be used. The resolution of screen dumps can be increased by reducing the print size of the original file while retaining the same number of pixels. You can also enlarge files by manipulating them in software such as PhotoShop. Your figures should be a minimum of 266 dpi when incorporated into your document.

\subsubsection{\LaTeX{} Overflow.}
\LaTeX{} users please beware: \LaTeX{} will sometimes put portions of the figure or table or an equation in the margin. If this happens, you need to scale the figure or table down, or reformat the equation. Check your log file! You must fix any overflow into the margin (that means no overfull boxes in \LaTeX{}). If you don't, the overflow text will simply be eliminated. \textbf{Nothing is permitted to intrude into the margins.}

\subsubsection{Using Color.}
AI Matters is published in color, so the use of color in figures is highly encouraged.

\subsubsection{Drawings.}
We suggest you use computer drawing software (such as Adobe Illustrator or, (if unavoidable), the drawing tools in Microsoft Word) to create your illustrations. Do not use Microsoft Publisher. These illustrations will look best if all line widths are uniform (half- to two-point in size), and you do not create labels over shaded areas. Shading should be 133 lines per inch if possible. Use Times Roman or Helvetica for all figure call-outs. \textbf{Do not use hairline width lines} --- be sure that the stroke width of all lines is at least .5 pt. Zero point lines will print on a laser printer, but will completely disappear on the high-resolution devices used by our printers.

\subsubsection{Photographs and Images.}
Photographs and other images should be in grayscale (color photographs will not reproduce well; for example, red tones will reproduce as black, yellow may turn to white, and so forth) and set to a minimum of 266 dpi. Do not prescreen images.

\subsubsection{Resizing Graphics.}
Resize your graphics \textbf{before} you include them with LaTeX. You may \textbf{not} use trim or clip options as part of your \textbackslash includgraphics command. Resize the media box of your PDF using a graphics program instead. 

\subsubsection{Fonts in Your Illustrations}
You must embed all fonts in your graphics before including them in your LaTeX document.




\section{Producing Reliable PDF\\Documents with \LaTeX{}}
Generally speaking, PDF files are platform independent and accessible to everyone. When creating a paper for a proceedings or publication in which many PDF documents must be merged and then printed on high-resolution PostScript RIPs, several requirements must be met that are not normally of concern. Thus to ensure that your paper will look like it does when printed on your own machine, you must take several precautions:
\begin{itemize}
\item Use type 1 fonts (not type 3 fonts)
\item Use only standard Helvet, Times, Nimbus, and CMR font packages (not fonts like F3 or fonts with tildes in the names or fonts---other than Computer Modern---that are created for specific point sizes, like Times\~{}19) or fonts with strange combinations of numbers and letters
\item Embed all fonts when producing the PDF
\item Do not use the [T1]{fontenc} package (install the CM super fonts package instead)
\end{itemize}

\subsection{Creating Output Using PDF\LaTeX{} Is Required}
By using the PDF\TeX{} program instead of straight \LaTeX{} or \TeX{}, you will probably avoid the type 3 font problem altogether (unless you use a package that calls for metafont). PDF\LaTeX{} enables you to create a PDF document directly from \LaTeX{} source. The one requirement of this software is that all your graphics and images must be available in a format that PDF\LaTeX{} understands (normally PDF).

PDF\LaTeX{}'s default is to create documents with type 1 fonts. If you find that it is not doing so in your case, it is likely that one or more fonts are missing from your system or are not in a path that is known to PDF\LaTeX{}.

\subsubsection{dvipdf Script}
Scripts such as dvipdf which ostensibly bypass the Postscript intermediary should not be used since they generally do not instruct dvips to use the config.pdf file.

\subsubsection{dvipdfm}
Do not use this dvi-PDF conversion package if your document contains graphics (and we recommend you avoid it even if your document does not contain graphics).

\subsection{Ghostscript}
\LaTeX{} users should not use GhostScript to create their PDFs.

\subsection{Graphics}
If you are still finding type 3 fonts in your PDF file, look at your graphics! \LaTeX{} users should check all their imported graphics files as well for font problems.

\section{Proofreading Your PDF}
Please check all the pages of your PDF file. Is the page size US Letter? Are there any type 3, Identity-H, or CID fonts? Are all the fonts embedded? Are there any areas where equations or figures run into the margins? Did you include all your figures? Did you follow mixed case capitalization rules for your title? Do any of the pages scroll slowly (because the graphics draw slowly on the page)? Are URLs underlined and in color? You will need to fix these common errors before submitting your file. 



\section{Additional Resources}
\LaTeX{} is a difficult program to master. If you've used that software, and this document didn't help or some items were not explained clearly, we recommend you read Michael Shell's excellent document (testflow doc.txt V1.0a 2002/08/13) about obtaining correct PS/PDF output on \LaTeX{} systems. (It was written for another purpose, but it has general application as well). It is available at www.ctan.org in the tex-archive.

\section{ Acknowledgments}
This document and the AI Matters style file were adapted from the AAAI formatting instructions and style file in 2015 by Eric Eaton.

We are especially grateful to Peter Patel Schneider for his work in implementing the original aaai.sty file, liberally using the ideas of other style hackers, including Barbara Beeton. We also acknowledge with thanks the work of George Ferguson for his guide to using the style and BibTeX files --- which has been incorporated into this document  --- and Hans Guesgen, who provided several timely modifications, as well as the many others who have, from time to time, sent in suggestions on improvements to the AAAI style. 

The preparation of the \LaTeX{} and Bib\TeX{} files that implement these instructions was supported by Schlumberger Palo Alto Research, AT\&T Bell Laboratories, Morgan Kaufmann Publishers, The Live Oak Press, LLC, and AAAI Press. Bibliography style changes were added by Sunil Issar. \verb+\+pubnote was added by J. Scott Penberthy. George Ferguson added support for printing the AAAI copyright slug. Additional changes to aaai.sty and aaai.bst have been made by the AAAI staff.

\bigskip
\noindent Thank you for reading these instructions carefully. We look forward to receiving your submissions to AI Matters!

\end{document}
